\chapter{Att kontrollera ett påstående}
\label{app:verifiera-tupler}

Det här materialet gör, som allt undervisningsmaterial, en rad
påståenden.  De flesta kan du pröva själv genom att köra programmen ---
det är hela poängen med att koden på bilderna är samma kod som körs.  Men
några påståenden går inte att pröva på det sättet:
\begin{enumerate}
  \item att en tupel i Python är oföränderlig, och att det är
    kommatecknet och inte parentesen som skapar den
    (\cref{sec:tupler,sec:uppackning});
  \item att \mintinline{python}{enumerate} lämnar tillbaka par, och inte
    en lista (\cref{sec:enumerate});
  \item att lösningen ska hållas enkel --- principen
    \foreignlanguage{english}{KISS} (\cref{sec:tupler});
  % Only the teacher's notes make such claims.
  \ifbool{ltnote}{%
  \item de didaktiska påståenden som anteckningarna i marginalen gör om
    hur människor lär sig.
  }{}%
\end{enumerate}
Hur vet man om sådant stämmer?  Samma fråga möter dig varje gång du läser
en lärobok, en webbsida eller ett svar från en språkmodell --- och varje
gång du själv ska skriva en rapport.
\Cref{tab:tupler-pastaenden} visar var vart och ett av materialets
påståenden är belagt.

\begin{table}[htbp]
  \centering
  \caption{Modulens påståenden, som frågor, och var svaret finns.}
  \label{tab:tupler-pastaenden}
  \begin{tabular}{@{}p{0.46\linewidth}p{0.44\linewidth}@{}}
    \toprule
    Fråga & Var det är belagt \\
    \midrule
    Är en tupel oföränderlig, och är det kommatecknet som skapar den?
      & Primärkälla: Pythons egen dokumentation, läst i fulltext
        2026-09-03; dessutom visar programmet i \cref{sec:tupler}
        felmeddelandet. \\
    Vad lämnar \mintinline{python}{enumerate} tillbaka?
      & Primärkälla: Pythons egen dokumentation; frågan ställs som
        övning så att studenten själv slår upp den, och programmet i
        \cref{sec:enumerate} visar typen och paren. \\
    Är enkelhet ett etablerat designvärde, och kommer
    \foreignlanguage{english}{KISS} från Kelly Johnson?
      & Redan belagt i föreläsningen \emph{Funktioner}, bilaga~C: ja på
        den första frågan, nej på den andra. \\
    % Rows about claims only the teacher's notes make.
    \ifbool{ltnote}{%
    Måste kontrast upplevas samtidigt för att en aspekt ska kunna
    urskiljas?
      & Primärkälla: Marton, \emph{Necessary Conditions of Learning},
        s.~45--47. \\
    }{}%
    \bottomrule
  \end{tabular}
\end{table}

Arbetsgången är densamma i alla kursens föreläsningar och beskrivs i
föreläsningen \emph{Algoritmiskt tänkande}, bilaga~A: hur ett påstående
görs till en fråga som kan besvaras, hur källan väljs efter frågan och
söks fram ämnesdrivet och tvåsidigt, med en motsökning efter
invändningar, hur fulltexten läses och spåret skrivs ner, och hur
slutsatsen dras med sina förbehåll.  Här tillämpas den på den här
föreläsningens påståenden i \cref{sec:tupler-pastaenden}.

\section{Materialets påståenden och var de är belagda}
\label{sec:tupler-pastaenden}

Inget av den här modulens påståenden krävde en egen litteratursökning.
Frågor om vad Python gör besvaras av Pythons egen
dokumentation\autocite{PythonDataStructures}, som är primärkällan: den
säger uttryckligen att tupler är oföränderliga och att deras element
kommer åt genom uppackning eller indexering.  Principen att välja den
enklaste lösning som duger (\foreignlanguage{english}{KISS}) namnges
här; det vetenskapliga underlaget för att enkelhet är ett etablerat
designvärde --- och för att tillskrivningen till Kelly Johnson inte går
att bekräfta --- finns i föreläsningen \emph{Funktioner}, bilaga~C.
% Marton backs claims only the teacher's notes make.
\ifbool{ltnote}{Att kontrast måste upplevas samtidigt för att en aspekt ska
kunna urskiljas är hämtat från Martons \emph{Necessary Conditions of
Learning}\autocite[45--47]{Marton2015}.  }{}Att återanvända ett belägg i
stället för att göra om sökningen är samma DRY-princip som kursen lär ut,
tillämpad på litteraturen.  Svaret på vart och ett av modulens påståenden
står sammanställt i \cref{tab:tupler-pastaenden}.

